% !TeX root = main.tex
% Обозначения, используемые в решениях.
\setlength{\emergencystretch}{2em}
\newcommand{\Z}{\mathbb{Z}}
\newcommand{\Q}{\mathbb{Q}}
\newcommand{\C}{\mathbb{C}}
\newcommand{\F}{\mathbb{F}}
\newcommand{\Rp}{\mathbb{R}_{>0}}
\DeclareMathOperator{\ord}{ord}
\DeclareMathOperator{\im}{Im}
\DeclareMathOperator{\lcm}{lcm}
\DeclareMathOperator{\sgn}{sgn}
\DeclareMathOperator{\diag}{diag}
\DeclareMathOperator{\GL}{GL}
\DeclareMathOperator{\SL}{SL}
\DeclareMathOperator{\Aut}{Aut}
\DeclareMathOperator{\Inn}{Inn}
\DeclareMathOperator{\Fix}{Fix}
\DeclareMathOperator{\Stab}{Stab}

% Номера задач берутся из исходного задания.
\newcommand{\task}[2]{\subsection*{\texorpdfstring{#1. #2}{Задача #1}}\addcontentsline{toc}{subsection}{\texorpdfstring{#1. #2}{Задача #1}}}
\newcommand{\editnote}[1]{\par\smallskip\noindent\textit{Примечание.}~#1\par\smallskip}
